\documentclass[UTF8,a4paper,12pt]{ctexart}

\usepackage{geometry}
\geometry{left=2.5cm,right=2.5cm,top=2.5cm,bottom=2.5cm}
\usepackage{graphicx}
\usepackage{booktabs}
\usepackage{amsmath}
\usepackage{hyperref}
\usepackage{listings}
\usepackage{xcolor}
\usepackage{tikz}
\usepackage{pgfplots}
\pgfplotsset{compat=1.18}
\usetikzlibrary{arrows.meta,positioning}

\lstset{
  basicstyle=\ttfamily\small,
  breaklines=true,
  frame=single,
  columns=fullflexible,
  showstringspaces=false
}

\title{基于开源多智能体框架的协作系统\\与工具调用 LoRA 微调研究}
\author{你的名字}
\date{\today}

\begin{document}

\maketitle

\begin{abstract}
本文围绕多智能体协作与工具调用格式稳定性展开。首先基于开源框架搭建包含任务分解、
执行与审查三个角色的多智能体系统；其次构造工具调用指令数据集，使用 QLoRA 对
Qwen2.5-7B-Instruct 进行轻量微调；最后在相同任务与工具集上比较单智能体、串行、
主管与辩论等协作模式，并分析微调前后 Agent 调用工具的 JSON 格式准确率与稳定性。
% TODO: 填入最终实验数字与结论。
\end{abstract}

\section{引言}
% TODO: 研究背景、问题定义、贡献列表。

\section{相关工作}
% TODO: 多智能体协作、ReAct、工具调用、LoRA/QLoRA。
多智能体系统通过角色分工与消息传递扩展单模型能力\cite{yao2023react}；
LoRA 以低秩适配降低微调成本\cite{hu2022lora}，QLoRA 进一步降低显存需求\cite{dettmers2023qlora}。

\section{系统设计}
\subsection{总体架构}
系统包含任务分解 Agent、执行 Agent 与审查/评估 Agent，三者通过共享状态与消息传递协作。
图~\ref{fig:architecture} 给出整体流程。

\begin{figure}[htbp]
  \centering
  \begin{tikzpicture}[
    node distance=1.5cm and 1.8cm,
    agent/.style={draw, rounded corners, minimum width=2.8cm,
      minimum height=1.0cm, align=center},
    >=Latex
  ]
    \node[agent] (planner) {Planner\\任务分解};
    \node[agent, right=of planner] (executor) {Executor\\工具调用};
    \node[agent, right=of executor] (reviewer) {Reviewer\\审查评估};

    \draw[->] (planner) -- (executor);
    \draw[->] (executor) -- (reviewer);
    \draw[->] (reviewer.north) -- ++(0,0.9) -| node[pos=0.25, above] {不通过，带 critique 重规划} (planner.north);
    \node[below=0.5cm of reviewer, align=center] {通过 $\rightarrow$ 输出};
  \end{tikzpicture}
  \caption{三角色多智能体协作流程}
  \label{fig:architecture}
\end{figure}

\subsection{角色与状态}
% TODO: 给出每个角色的 prompt 要点、输入输出契约、停止条件。
共享状态包括任务、计划、执行结果、审查结果、轮数与 trace。

\subsection{工具层}
% TODO: 工具注册表、JSON Schema、错误处理、重试。
工具调用统一为：
\begin{lstlisting}
{"name": "calculator", "arguments": {"expression": "(23*17+9)/5"}}
\end{lstlisting}

\section{工具调用数据与微调}
\subsection{数据构造}
% TODO: 数据规模、工具覆盖、硬负例、按工具划分训练/测试集。

\subsection{QLoRA 配置}
% TODO: rank、alpha、target modules、学习率、epoch、显存与耗时。

\section{实验设置}
\subsection{任务与指标}
格式指标包括 JSON 可解析率、Schema 合规率、工具名准确率、参数完全匹配率；
系统指标包括任务成功率、平均重试次数、延迟与 token 成本。

\subsection{对照与协作模式}
% TODO: base / few-shot / 约束解码 / LoRA / LoRA+约束解码。
% TODO: 单 Agent、串行、主管、辩论四种拓扑。

\section{结果与分析}
% TODO: 表格与图。建议至少包含：
% 1) 微调前后格式指标对比；
% 2) 未见工具 OOD 结果；
% 3) 四种协作模式的成功率-成本权衡；
% 4) 失败案例分类。

\section{讨论与局限性}
% TODO: LoRA 与约束解码的边界、错误传播、成本、评测偏差。

\section{结论}
% TODO: 总结主要发现。

\bibliographystyle{plain}
\bibliography{refs}

\end{document}
